\chapter{Percutaneous Endoscopic Gastrostomy}

\section*{Overview}
Percutaneous endoscopic gastrostomy (PEG) is a procedure that places a feeding tube directly into the stomach through the abdominal wall using endoscopic guidance, for long-term enteral nutrition.

\section*{Alternative Names}
PEG tube placement; gastrostomy tube; feeding tube; percutaneous gastrostomy

\section*{Principle}
An endoscope is passed into the stomach, and the stomach wall is pushed against the abdominal wall. A needle is advanced percutaneously into the stomach, a guidewire is introduced and pulled through the mouth, and a feeding tube is drawn back through the abdominal wall into the stomach.

\section*{History}
Percutaneous endoscopic gastrostomy was introduced by Gauderer, Ponsky, and Izant in 1980 and rapidly replaced surgical gastrostomy for most patients needing long-term tube feeding.

\section*{Why It Is Done}
Performed for patients who cannot maintain adequate oral nutrition, such as those with dysphagia from stroke, neurologic disease, head and neck cancer, or severe malnutrition.

\section*{Is This a Primary or Secondary Test or Procedure}
Primary long-term enteral access procedure.

\section*{How It Is Performed}
Performed under conscious sedation or anesthesia. An endoscope confirms position and transillumination guides the puncture site. After prepping, the stomach is entered percutaneously, the tube is placed, and its position is confirmed endoscopically. The tube is secured at the skin.

\section*{Preparation}
NPO before the procedure, IV antibiotics, and assessment of coagulation. The patient's head is elevated, and the puncture site is chosen below the xiphoid.

\section*{Contraindications and Precautions}
Total gastrectomy, severe coagulopathy, and severe ascites are contraindications. Precaution: prior abdominal surgery with interposed bowel may preclude safe placement.

\section*{Risks}
Bleeding, infection at the site, peritonitis, perforation, tube migration or blockage, aspiration, and buried bumper syndrome.

\section*{Aftercare and Recovery}
Feeds begin within hours, and the site is kept clean and dry. The tract matures over 2-4 weeks, after which the tube can be replaced if needed.

\section*{Outcome and Follow-up}
Successful tube placement is confirmed endoscopically; feeds are tolerated and weight stabilizes or improves.

\section*{Expected Results}
A functioning tube delivering nutrition with a healed, non-infected site.

\section*{Follow-up}
Ongoing nutrition monitoring, tube care, and replacement as the tube wears or if the patient regains oral intake.

\section*{References}
\begin{enumerate}
  \item MedlinePlus. PEG tube. U.S. National Library of Medicine; 2025.
  \item Current Medical Diagnosis and Treatment. McGraw-Hill; 2025.
\end{enumerate}

\clearpage
